\chapter{Ochrona prywatności przez selektywne ujawnianie danych}
\label{ch:prywatnosc}

System podstawowy zakotwicza w~łańcuchu bloków skrót całej treści dyplomu.
Rozdział ten pokazuje, jak --- bez żadnej zmiany w~kontrakcie --- zastąpić ten
monolityczny skrót strukturą zobowiązań kryptograficznych, dzięki której posiadacz
dyplomu może ujawnić weryfikatorowi jedynie wybrane pola, zachowując pozostałe
w~tajemnicy.

\section{Problem: minimalizacja danych}
\label{sec:priv-problem}

W~systemie podstawowym skrót dokumentu powstaje przez kanonizację i~zahaszowanie
\emph{całego} ładunku dyplomu: \code{docHash = keccak256(canonicalize(payload))}.
Aby zweryfikować dyplom, posiadacz musi przekazać weryfikatorowi pełny ładunek ---
inaczej weryfikator nie odtworzy skrótu. Oznacza to, że dowód posiadania stopnia
naukowego wymusza ujawnienie wszystkich danych: imienia i~nazwiska, numeru dyplomu,
daty wydania i~szczegółów kierunku, nawet gdy relacja weryfikacyjna wymaga
potwierdzenia tylko jednego z~nich.

Narusza to zasadę minimalizacji danych --- fundament nowoczesnych regulacji ochrony
danych osobowych --- zgodnie z~którą przetwarzać należy jedynie dane niezbędne do
celu. Pracodawca weryfikujący sam fakt ukończenia studiów nie powinien być zmuszany
do poznania numeru dyplomu kandydata. Celem tego rozdziału jest umożliwienie
weryfikacji \emph{pojedynczych pól} bez ujawniania reszty, przy zachowaniu tej samej
gwarancji zakotwiczenia w~łańcuchu.

\section{Konstrukcja zobowiązań pól}
\label{sec:priv-zobowiazania}

Ładunek dzielimy na cztery ujawnialne pola: \code{student}, \code{degree},
\code{issuedAt} oraz \code{diplomaNumber}. Dla każdego z~nich obliczamy osobne
\emph{zobowiązanie ukrywające} (\emph{hiding commitment}), łącząc kanoniczną postać
wartości pola z~losową solą:
\[
  c_i = \operatorname{keccak256}\bigl(\operatorname{canonicalize}(\text{pole}_i)
        \,\Vert\, \text{sól}_i\bigr),
  \qquad i \in \{\text{student}, \text{degree}, \text{issuedAt},
        \text{diplomaNumber}\}.
\]
Każda sól to niezależna, losowa wartość 256-bitowa generowana przez
\code{crypto.getRandomValues}. Skrót dokumentu odtwarzamy następnie \emph{ze
zobowiązań}, a~nie z~jawnych wartości:
\[
  \code{docHash} = \operatorname{keccak256}
    \bigl(c_{\text{student}} \Vert c_{\text{degree}}
    \Vert c_{\text{issuedAt}} \Vert c_{\text{diplomaNumber}}\bigr).
\]
Liść drzewa Merkle pozostaje związany z~kontekstem wdrożenia dokładnie tak, jak
w~systemie podstawowym:
\[
  \code{leaf} = \operatorname{keccak256}
    \bigl(\code{registry} \Vert \code{chainId} \Vert \code{issuer}
    \Vert \code{batchId} \Vert \code{docHash}\bigr).
\]
Ponieważ \code{docHash} jest deterministyczną funkcją czterech zobowiązań, korzeń
partii publikowany w~kontrakcie nie zmienia formatu --- kontrakt nie odróżnia
dyplomu prywatnego od standardowego.

\section{Koperta prywatna}
\label{sec:priv-koperta}

Prywatna koperta dyplomu (\code{PrivateDiplomaEnvelope}) zawiera trzy części:

\begin{itemize}
  \item \textbf{zobowiązania} --- pełny komplet czterech wartości $c_i$, potrzebny
        do odtworzenia \code{docHash} niezależnie od tego, które pola ujawniono;
  \item \textbf{pola ujawnione} --- dla każdego \emph{wybranego} pola jego jawna
        wartość wraz z~odpowiadającą jej solą; pola nieujawnione są w~kopercie
        nieobecne;
  \item \textbf{dowód Merkle} --- ścieżka od liścia do korzenia partii, identyczna
        jak w~kopercie standardowej.
\end{itemize}

Strukturę koperty i~drogę od pól do korzenia partii ilustruje
rys.~\ref{fig:koperta}.

\begin{figure}[H]
\centering
\begin{tikzpicture}[node distance=4mm and 7mm]
  \node[blok] (p1) {\code{student}\\\footnotesize + sól$_1$};
  \node[blok, right=of p1] (p2) {\code{degree}\\\footnotesize + sól$_2$};
  \node[blok, right=of p2] (p3) {\code{issuedAt}\\\footnotesize + sól$_3$};
  \node[blok, right=of p3] (p4) {\code{diplomaNumber}\\\footnotesize + sól$_4$};
  \node[blokakcent, below=7mm of p1] (c1) {$c_1$};
  \node[blokakcent, below=7mm of p2] (c2) {$c_2$};
  \node[blokakcent, below=7mm of p3] (c3) {$c_3$};
  \node[blokakcent, below=7mm of p4] (c4) {$c_4$};
  \node[blokszary, below=8mm of $(c2)!0.5!(c3)$, minimum width=50mm]
    (dochash) {\code{docHash} $= H(c_1 \Vert c_2 \Vert c_3 \Vert c_4)$};
  \node[blokszary, below=5mm of dochash, minimum width=50mm]
    (leaf) {liść $= H(\code{registry} \Vert \code{chainId} \Vert \code{issuer}
            \Vert \code{batchId} \Vert \code{docHash})$};
  \node[blokszary, below=5mm of leaf] (root) {korzeń partii (on-chain)};
  \foreach \i in {1,2,3,4}
    \draw[strz] (p\i) -- node[etyk, right=0.5mm] {$H$} (c\i);
  \draw[strz] (c1.south) -- (dochash.north west);
  \draw[strz] (c2) -- (dochash);
  \draw[strz] (c3) -- (dochash);
  \draw[strz] (c4.south) -- (dochash.north east);
  \draw[strz] (dochash) -- (leaf);
  \draw[strz] (leaf) -- node[etyk, right=0.5mm] {dowód Merkle} (root);
\end{tikzpicture}
\caption[Struktura koperty prywatnej]{Koperta prywatna: każde pole wraz z~tajną solą daje zobowiązanie
$c_i$; skrót dokumentu powstaje z~zobowiązań, nie z~jawnych wartości. Pole
ujawnione = wartość + sól; pole ukryte = samo $c_i$.}
\label{fig:koperta}
\end{figure}

Sól pełni tu podwójną rolę. Dla pola \emph{ujawnionego} jest przekazywana
weryfikatorowi, aby mógł odtworzyć i~sprawdzić zobowiązanie. Dla pola
\emph{ukrytego} pozostaje tajemnicą posiadacza --- to ona uniemożliwia odgadnięcie
wartości pola z~samego zobowiązania.

\section{Protokół weryfikacji selektywnej}
\label{sec:priv-protokol}

Weryfikator dysponujący prywatną kopertą wykonuje trzy kroki:

\begin{enumerate}
  \item \textbf{Sprawdzenie ujawnionych pól.} Dla każdego pola obecnego w~sekcji
        ujawnionej weryfikator przelicza $c_i' = \operatorname{keccak256}
        (\operatorname{canonicalize}(\text{wartość}) \Vert \text{sól})$ i~porównuje
        z~zobowiązaniem $c_i$ z~koperty. Niezgodność któregokolwiek pola oznacza
        odrzucenie.
  \item \textbf{Odtworzenie skrótu dokumentu.} Ze wszystkich czterech zobowiązań
        (ujawnionych i~ukrytych) weryfikator liczy \code{docHash}, a~z~niego liść.
  \item \textbf{Weryfikacja dowodu Merkle.} Odtworzony liść wraz z~dowodem daje
        korzeń, który weryfikator porównuje z~korzeniem odczytanym z~kontraktu dla
        pary \code{(issuer, batchId)}.
\end{enumerate}

Pole ukryte przechodzi przez krok~2 jako nieprzejrzyste zobowiązanie: jego wartość
nie jest weryfikatorowi znana, ale jego wkład do \code{docHash} --- a~więc i~do
liścia --- jest w~pełni uwzględniony. Weryfikator zyskuje pewność, że ujawnione pola
należą do dyplomu zakotwiczonego on-chain, nie dowiadując się niczego o~polach
ukrytych.

\section{Implementacja}
\label{sec:priv-implementacja}

Cała logika prywatności jest \emph{poza łańcuchem}. Kontrakt \code{DiplomaRegistry}
pozostaje niezmieniony --- publikuje korzenie przez ten sam \code{issueBatchRoot},
a~jego koszty gazu są identyczne jak w~rozdziale~\ref{ch:l2}. Zmiany ograniczają się
do warstwy klienckiej: pakiet \code{verifier-core} zyskuje funkcje
\code{generateFieldSalts}, \code{computeFieldCommitments}, \code{computePrivateDocHash}
i~\code{verifyDisclosedFields} (pokryte 17~testami jednostkowymi), SDK udostępnia
weryfikację selektywną, a~aplikacja webowa --- trasę \code{/present}, na której
posiadacz zaznacza pola do ujawnienia i~generuje prywatną kopertę.

Fakt, że tak istotna funkcja realizowana jest bez ani jednej zmiany w~kontrakcie,
jest sam w~sobie wynikiem: pokazuje, że model zakotwiczenia korzenia Merkle jest na
tyle ogólny, iż selektywne ujawnianie mieści się w~nim jako czysto klienckie
rozszerzenie.

\section{Analiza bezpieczeństwa}
\label{sec:priv-bezpieczenstwo}

Bezpieczeństwo konstrukcji opiera się na dwóch własnościach funkcji
\code{keccak256} traktowanej jako zobowiązanie.

\paragraph{Wiązanie (binding).} Odporność \code{keccak256} na kolizje sprawia, że
posiadacz nie może otworzyć jednego zobowiązania na dwie różne wartości pola.
Uniemożliwia to przedstawienie tego samego dyplomu z~podmienioną, lecz „poprawnie
weryfikującą się'' zawartością.

\paragraph{Ukrywanie (hiding).} Dla pola ukrytego napastnik widzi jedynie
$c_i = \operatorname{keccak256}(\text{wartość} \Vert \text{sól})$. Atak wyczerpujący
wymagałby zgadnięcia \emph{zarówno} wartości, \emph{jak i} 256-bitowej soli. Ponieważ
sól pochodzi z~kryptograficznie bezpiecznego generatora i~pozostaje tajna, przestrzeń
poszukiwań jest rzędu $2^{256}$ --- praktycznie nieprzeszukiwalna, nawet gdy sama
wartość pola należy do małej domeny (np. \code{degree.level} przyjmuje kilka
wartości). Odsala to typowy zarzut wobec zobowiązań nad małymi domenami: bez
tajnej soli o~wysokiej entropii atak słownikowy byłby trywialny, lecz tutaj sól tę
lukę zamyka.

\paragraph{Odporność na powtórzenie (replay).} Związanie liścia z~czwórką
\code{(registry, chainId, issuer, batchId)} uniemożliwia przeniesienie ważnego
dyplomu do innego kontraktu, sieci, wydawcy lub partii.

\paragraph{Ograniczenie.} Konstrukcja ukrywa \emph{wartości} pól, lecz nie ukrywa
\emph{struktury} koperty --- liczby pól ani faktu, że dyplom jest prywatny. Nie
zapewnia też dowodów o~ujawnionej wartości typu „stopień jest co najmniej
licencjatem'' bez podania konkretnej wartości. Takie własności wymagałyby dowodów
z~wiedzą zerową (ZK), które --- zgodnie z~przyjętym zakresem pracy --- pozostają poza
jej ramami; obecna konstrukcja realizuje selektywne ujawnianie pełnych pól, co
odpowiada założonemu wymaganiu minimalizacji danych. Funkcjonalnie odpowiada to
selektywnemu ujawnianiu znanemu z~poświadczeń weryfikowalnych
W3C~\cite{w3c-vc}, przy czym tam realizuje się je specjalizowanymi schematami
podpisów, a~tu --- wyłącznie funkcją skrótu i~solą, kosztem ujawnienia
struktury koperty.

\section{Zgodność wstecz}
\label{sec:priv-zgodnosc}

Koperta standardowa działa bez zmian --- weryfikator, który nie zna rozszerzenia
prywatnego, przetwarza ją tak jak dotąd. Na poziomie partii przyjmujemy regułę
rozłączności: dana partia jest w~całości standardowa albo w~całości prywatna, co
usuwa dwuznaczność formatu \code{docHash} przy odtwarzaniu liścia i~pozwala obu
trybom współistnieć w~jednym kontrakcie.
